\bookchapter{第二十一章：移動網絡與資訊安全深探}

\chapterintro{我哋每日用手機上網、支付、傾偈，感覺方便得好。但你有冇諗過，你嘅傾偈記錄、支付密碼、相片……呢啲資訊喺網上“裸跑”嘅時候多嗎？講真，大部分人根本冇諗過呢個問題，但一旦出事——比如帳號被盜、銀行卡被碌——嗰陣時再後悔就遲啦。呢一章我哋就好好傾下，點樣等你嘅數字生活唔咁“透明”。}

\booksection{一、VPN：網絡入面嘅“隱形斗篷”}

VPN，全稱虛擬專用網絡（Virtual Private Network）。你可以將佢想像成喺你屋企同你要去嘅網站之間，建起咗一條專屬嘅、加密嘅地下隧道。你喺隧道入面行來行去，外面嘅人睇唔到你手入面揸住乜嘢、要去邊。

\begin{itemize}
  \item \textbf{普通上網}：你嘅手機公共網絡目標網站。中間經過嘅路由器、交換機，隨便邊個環節都可能有人"偷睇"你發送嘅數據。尤其係喺咖啡廳、機場嗰種公共Wi-Fi上，講難聽啲，簡直係"裸跑"。
  \item \textbf{使用VPN}：你嘅手機VPN加密隧道VPN伺服器目標網站。數據被加密封裝，"偷窺者"淨係睇到一堆亂碼。另外，目標網站睇到嘅係VPN伺服器嘅IP地址，而唔係你嘅真實IP。
\end{itemize}

噉VPN到底能做乜？簡單講：第一，\textbf{保護你喺公共Wi-Fi上嘅隱私}；第二，\textbf{繞過地域限制}（比如睇某啲淨係喺特定國家上線嘅內容）；第三，\textbf{防止你嘅網絡營辦商（ISP）記錄你嘅瀏覽習慣}。不過我要提醒一句：\textbf{唔好用免費VPN！}免費嘅背後往往意味住你嘅數據就係佢哋嘅“商品”。我有個朋友用咗某個免費VPN，結果手機莫名其妙開始彈廣告，後來發現嗰個VPN將佢嘅瀏覽數據打包賣畀廣告商啦。所以揀VPN，要麼揀付費嘅大品牌（如NordVPN、ExpressVPN），要麼乾脆唔好用。

另外，VPN唔係萬能嘅。佢淨係可以保護“傳輸中”嘅數據，如果目標網站本身唔安全、或者你嘅手機上已經中咗惡意軟件，VPN都幫唔到你。

\booksection{二、公共Wi-Fi嘅“甜蜜陷阱”——比你想像嘅危險}

你可能覺得，唔就係連個免費Wi-Fi嘛，能有多大風險？我同你講個真事：安全研究人員曾經做過實驗，喺機場擺咗一個叫“Starbucks\_Free”嘅Wi-Fi熱點，唔夠一個鐘就有上百人連接。而實際上呢個熱點係研究人員自己搭嘅，所有經過呢個熱點嘅未加密流量——比如你訪問嘅HTTP網頁內容、未啟用TLS嘅App通信——佢哋確實能截獲同埋查看。

但要注意一個關鍵前提：今日你日常用嘅微信、網銀、支付寶，絕大多數都已經行咗HTTPS加密。講白啲，攻擊者即使截獲咗數據包，睇到嘅都淨係一堆亂碼，\textbf{無法直接讀出你嘅傾偈內容或支付密碼}。不過佢哋依然可以知道你訪問咗邊啲網站（域名），如果某啲老舊App或網頁仲喺度用HTTP明文傳輸，噉內容就會完全暴露。呢種攻擊叫“中間人攻擊”（Man-in-the-Middle），原理係攻擊者喺你同真正嘅Wi-Fi路由器之間“插咗一腳”。更可怕嘅係，好多人嘅手機設定咗“自動連接已知Wi-Fi”，攻擊者只要將自己嘅熱點名設成你之前連過嘅某個Wi-Fi（比如“ChinaNet”或“CMCC”），你嘅手機就會自動連上去。

噉點樣防？幾條實用建議：

\begin{itemize}
  \item \textbf{認準HTTPS}：地址欄開頭係https://，多出嘅"s"代表Secure。訪問任何需要輸密碼嘅網站時，如果睇到嘅係http://（冇s），即刻停低唔好輸！而家主流瀏覽器都會喺非HTTPS網站前顯示"唔安全"嘅警告，睇到呢個標記就要警惕。
  \item \textbf{關掉"自動連接"}：唔論係iPhone定係安卓，都喺Wi-Fi設定入面將"自動加入"或"自動連接"關掉。寧可每次手動點一下，都唔好等手機自己判斷。
  \item \textbf{用手機流量做敏感操作}：轉賬、登入網銀、輸入身份證號……呢啲操作老老實實用自己嘅4G/5G流量，唔好慳嗰啲流量錢。
  \item \textbf{忘記唔常用嘅Wi-Fi}：出差返嚟、旅遊返嚟，將之前臨時連過嘅酒店Wi-Fi、機場Wi-Fi都"忘記"掉。噉下次手機唔會自動連到同名假熱點。
\end{itemize}

\booksection{三、HTTPS係點樣保護你嘅？——SSL/TLS握手}

啱啱提到咗HTTPS，噉佢到底係點樣工作嘅呢？每次你訪問一個HTTPS網站，瀏覽器同伺服器之間會進行一次“SSL/TLS握手”，成個過程大概就係噉：

\begin{enumerate}
  \item 你嘅瀏覽器話："嘿，伺服器，我想同你安全通信，呢個係我嘅加密能力清單（支援嘅加密算法等）。"
  \item 伺服器回覆："好嘅，我從你嘅清單入面揀咗呢套加密方案，呢個係我嘅數字證書，入面麵包含我嘅公鑰，由某某CA機構簽名認證。"
  \item 瀏覽器驗證證書嘅真假：係唔係由可信嘅CA簽發？有冇過期？域名對唔對得上？
  \item 驗證通過後，雙方用公鑰加密協商出一個"會話密鑰"，之後嘅通信全用呢個密鑰加密。
\end{enumerate}

呢個就似你去銀行辦事：你先遞上身份證（證書），櫃員核實你嘅身份，然後你哋兩個約定一個暗號，之後嘅對話淨係得你哋兩個聽得明。即使有人喺旁邊偷聽，都淨係聽到加密後嘅亂碼。

但注意：HTTPS淨係保證“傳輸過程”係加密嘅，唔保證目標網站係好人。釣魚網站都可以用HTTPS——嗰個小鎖圖示淨係說明你連接嘅係“某個網站”，唔代表嗰個網站就係你想去嘅嗰個。所以都係要自己核對網址係咪正確。

\booksection{四、DNS安全：你嘅“網絡電話簿”安全嗎？}

DNS（域名系統）就似互聯網嘅電話簿：你輸入“baidu.com”，DNS幫你查到佢對應邊個IP地址（比如110.242.68.66），然後瀏覽器先可以連過去。問題係，傳統嘅DNS查詢係明文嘅、冇加密嘅！等於話，你嘅ISP、Wi-Fi提供者、或者任何中間節點，都能睇到你訪問過邊啲網站。

更惡心嘅係“DNS劫持”：攻擊者喺DNS環節做手腳，將“baidu.com”指向一個假嘅釣魚網站。你明明輸入嘅係正確網址，卻被帶到咗一個一模一樣但實際係假嘅登入頁面，一輸入帳號密碼就玩完啦。

好喺而家有咗加密DNS方案：

\begin{itemize}
  \item \textbf{DNS over HTTPS（DoH）}：將DNS查詢收埋喺HTTPS流量入面，混喺正常嘅網頁瀏覽流量中，根本分唔出邊啲係DNS查詢。
  \item \textbf{DNS over TLS（DoT）}：用TLS加密協議專門保護DNS查詢，埠獨立，網管能知道你用咗DoT但睇唔到內容。
\end{itemize}

我個人推薦用Cloudflare嘅1.1.1.1或者Google嘅8.8.8.8作為DNS伺服器，佢哋都支援加密DNS。喺安卓嘅“私人DNS”設定入面直接填“one.one.one.one”或者“dns.google”就得，超簡單。iPhone稍微麻煩啲，需要裝個描述檔或者通過App設定，但都值得搞下。

\booksection{五、“無痕瀏覽”真係能隱身嗎？}

先講結論：\textbf{無痕模式係你屋企窗簾，唔係隱身斗篷。}

佢做咗乜：唔喺本地留低痕跡——關閉無痕視窗後，瀏覽器唔保存你嘅瀏覽歷史、Cookies、表單填寫記錄、搜索記錄。借人哋電腦查嘢、或者買禮物唔想被家人從瀏覽記錄發現，用無痕模式係合適嘅。

佢做唔到乜：

\begin{itemize}
  \item 你訪問嘅網站\textbf{照樣能睇到你嚟咗}，你嘅IP地址、設備資訊、瀏覽行為，網站端一清二楚。
  \item 你嘅網絡營辦商\textbf{照樣能睇到你訪問咗邊啲網站}，淨係睇唔到具體內容（如果用咗HTTPS）。
  \item 公司/學校嘅網絡管理員\textbf{一樣能監控}你嘅上網行為。
  \item 你下載嘅檔案、添加嘅書籤會\textbf{保留落嚟}，唔會隨住關閉視窗而消失。
\end{itemize}

所以記住：無痕模式係“本地擦除器”，唔係“網絡隱身衣”。如果你真係需要隱藏上網痕跡，就要用VPN+Tor瀏覽器呢種組合拳，不過對於日常使用嚟講，瞭解呢啲就夠用啦。

\booksection{六、瀏覽器隱私：唔好等瀏覽器將你賣咗}

你可能冇意識到，瀏覽器本身就係一個巨大嘅資訊收集器。Cookies、瀏覽歷史、快取、自動填充、位置權限……每樣都可能洩露你嘅隱私。幾個建議：

\begin{itemize}
  \item \textbf{定期清理瀏覽數據}：每隔一段時間，去設定入面將Cookies、快取、歷史記錄清一清。唔係為咗防止"被追蹤"——嗰個靠清Cookies已經唔太管用啦——而係為咗防止手機越來越卡。瀏覽器快取能佔好幾個G嘅空間。
  \item \textbf{關掉"請勿追蹤"嘅幻想}：瀏覽器入面嘅"發送'請勿追蹤'請求"選項，本質上係"請"網站唔好追蹤你。但網站聽唔聽完全係另一回事，大部分網站直接忽略。唔好指望佢。
  \item \textbf{善用隱私保護瀏覽器}：Firefox Focus、DuckDuckGo瀏覽器、Brave呢啲瀏覽器預設攔截追蹤器同埋廣告，比Chrome同Safari嘅預設設定要乾淨得多。
  \item \textbf{搜索引擎都重要}：用DuckDuckGo代替百度/Google搜索，佢唔記錄你嘅搜索歷史，唔建立用戶畫像。講真搜索結果確實唔如Google精準，但日常查資料夠用啦。
  \item \textbf{檢查App權限}：去手機設定入面睇下你嘅瀏覽器被授予咗乜嘢權限。佢真係需要訪問你嘅位置、咪、通訊錄嗎？大概率唔需要。
\end{itemize}

\booksection{七、實用安全清單：保護自己從而家開始}

咁樣講咗咁多，最後畀你一個簡單粗暴嘅安全清單，花半個鐘搞掂，你嘅手機安全性會有質嘅飛躍：

\begin{enumerate}
  \item \textbf{關掉Wi-Fi自動連接}：iPhone喺設定-無線局域網入面，安卓喺Wi-Fi設定入面。
  \item \textbf{換一個加密DNS}：安卓直接設定私人DNS為"dns.google"，iPhone用App或描述檔設定1.1.1.1嘅DoH。
  \item \textbf{開啟雙因素認證（2FA）}：微信、支付寶、Apple ID、Google帳號，能開2FA嘅都開。短訊驗證碼已經係最低級別嘅2FA啦，能用TOTP（如Google Authenticator、Microsoft Authenticator）更好。
  \item \textbf{查一遍App權限}：將每個App嘅權限過一遍，特別係位置、通訊錄、短訊、咪、相機呢五項。你會發現有啲電筒App居然要你嘅通訊錄——即刻卸載。
  \item \textbf{定期清理瀏覽器數據}：設定個日曆提醒，每月一次，清快取、Cookies。
  \item \textbf{唔好用公共充電樁嘅數據線}：USB口既能充電都能傳數據，惡意充電樁（叫"Juice Jacking"）可以趁你充電時偷數據。出街帶自己嘅充電頭同埋線。
\end{enumerate}

\begin{tcbpractice}{而家就做：畀你嘅手機做一次"安全體檢"}
\begin{enumerate}
  \item \textbf{DNS升級}：安卓：設定 > 連接與共享 > 私人DNS > 填入"one.one.one.one"或"dns.google"。iPhone：用Safari搜索"1.1.1.1 DNS profile"安裝描述檔，或下載Cloudflare嘅1.1.1.1 App。
  \item \textbf{Wi-Fi審計}：打開Wi-Fi列表，將唔識嘅、唔常用嘅、名稱可疑嘅熱點都"忘記"掉。
  \item \textbf{權限大掃除}：進入設定 > 隱私/權限管理，逐個App審查權限。標記出嗰啲權限要求唔合理（比如計算機要通訊錄）嘅App，考慮卸載或用替代品。
  \item \textbf{檢查帳號安全}：打開Apple ID或Google帳號嘅安全頁面，睇下有冇不明設備登入，確認2FA已開啟，檢查恢復聯絡人係咪可信。
\end{enumerate}

\end{tcbpractice}
